% 三版共用的字体：正文与导航、单色符号、数学、行内等宽与代码块。
% 由各版布局在 \geometry 之后 \input；字号、行距等参数由各版布局设置。
% 本机字体须可供 XeLaTeX 访问，不静默替换。

% 正文与导航字体：使用与已确认样张一致的 Latin Modern Roman 10 pt 静态字形。
\setmainfont{lmroman10}[
  Extension=.otf,UprightFont=*-regular,BoldFont=*-bold,
  ItalicFont=*-italic,BoldItalicFont=*-bolditalic
]
\setsansfont{PingFang SC}[BoldFont={PingFang SC Semibold}]
\setCJKmainfont{Songti SC}[BoldFont={Songti SC Bold},ItalicFont={Songti SC}]
\setCJKsansfont{PingFang SC}[BoldFont={PingFang SC Semibold}]
% 宋体正文的双引号改用当前西文正文字形（用户确认），输入仍为 Unicode “”。
% 跟随 NFSS 字体族选择：粗体/斜体保留规则，苹方与 Maple（含自定义代码字体）恢复原分类。
% 保存 xeCJK 已分配的类号，不硬编码；半角左右标点类取自西文圆括号。
% 不在 selectfont 中反复调用 xeCJKDeclareCharClass：它会全局追加字符范围表。
% 这里只作局部 XeTeXcharclass 赋值，嵌套字体分组结束时自然恢复外层分类。
\makeatletter
\chardef\BookQuoteFullLeft=\XeTeXcharclass"201C
\chardef\BookQuoteFullRight=\XeTeXcharclass"201D
\chardef\BookQuoteHalfLeft=\XeTeXcharclass`\(
\chardef\BookQuoteHalfRight=\XeTeXcharclass`\)
\newcommand{\BookSelectQuoteClass}{%
  \ifx\f@family\rmdefault
    \XeTeXcharclass"201C=\BookQuoteHalfLeft
    \XeTeXcharclass"201D=\BookQuoteHalfRight
  \else
    \XeTeXcharclass"201C=\BookQuoteFullLeft
    \XeTeXcharclass"201D=\BookQuoteFullRight
  \fi}
\AddToHook{selectfont}[book/quotes]{\BookSelectQuoteClass}
\makeatother
% 单色符号使用项目内的静态字体；标题中也保持常规直立字重。
% 只映射原稿实际使用的两个字符，不更改正文或全局字符分类。
\newfontfamily\BookSymbolFont{NotoEmoji-Regular.ttf}[
  Path=preamble/fonts/,
  BoldFont=NotoEmoji-Regular.ttf,
  ItalicFont=NotoEmoji-Regular.ttf,
  BoldItalicFont=NotoEmoji-Regular.ttf
]
% xeCJK 默认借用当前字体空格；Noto Emoji 的空格宽达 1.27em。
% 仅在符号局部分组内固定中西文间距，避免中文后的符号被推远。
\DeclareRobustCommand{\BookStar}{{%
  \xeCJKsetup{CJKecglue={\hskip .25em\relax}}%
  \BookSymbolFont\char"2B50\relax}}
\DeclareRobustCommand{\BookFlag}{{%
  \xeCJKsetup{CJKecglue={\hskip .25em\relax}}%
  \BookSymbolFont\char"1F6A9\relax}}
\newunicodechar{⭐}{\BookStar}
\newunicodechar{🚩}{\BookFlag}
% 书签保留原 Unicode，避免把字体命令写入 PDF 字符串。
\pdfstringdefDisableCommands{%
  \def⭐{\string⭐}%
  \def🚩{\string🚩}%
}
\setmathfont{Latin Modern Math}
% 直立数学字母与算子均使用数学字体；\text 仍使用正文字体。
% 显式外层分组兼容原稿的 ^\mathrm{mid}、_\mathrm{mid} 写法。
% 直接使用 mathrm=sym 时，unicode-math 的内部组命令不能充当上下标组。
% 在 unicode-math 完成文档起始初始化后设置，避免被其默认定义覆盖。
\AddToHook{begindocument/end}{\renewcommand{\mathrm}[1]{{\symup{#1}}}}
\setoperatorfont{\symup}
% 保留 \boldsymbol 对完整表达式（含算子和定界符）的语义。
% 数学字体没有独立粗体文件，bold 版本采用轻度合成加粗。
\setmathfont{Latin Modern Math}[version=bold,FakeBold=1.2]
% 保留文本类数学字母（如 \mathbf、\mathit）的 bold 版本字体映射。
% \mathrm 已由上方定义改用数学字形，不走此文本映射。
\setboldmathrm{lmroman10}[
  Extension=.otf,UprightFont=*-bold,BoldFont=*-bold,
  ItalicFont=*-bolditalic,BoldItalicFont=*-bolditalic
]

% 行内等宽文字使用 Maple，字号随所在正文；与代码块采用相同字形特性。
\setmonofont{Maple Mono NF CN}[
  Ligatures={NoCommon,NoContextual,NoDiscretionary,NoHistoric},
  RawFeature={-calt;-tlig;+cv01;+cv32;+cv33;+cv34;+cv35;+cv36;+cv37}
]
\setCJKmonofont{Maple Mono NF CN}[
  Ligatures={NoCommon,NoContextual,NoDiscretionary,NoHistoric},
  RawFeature={-calt;-tlig;+cv01;+cv32;+cv33;+cv34;+cv35;+cv36;+cv37}
]

% 只供 minted 使用（fontfamily=bookcode）；不改变 \texttt 或代码标题栏。
% 禁用标准、上下文、可选、历史连字，以及 calt/TeX 字符替换。
% cv01 去除符号间隙；cv32--cv37 将手写斜体改为常规斜体（仍保留倾斜）。
\newfontfamily\BookCodeLatin{Maple Mono NF CN}[
  Ligatures={NoCommon,NoContextual,NoDiscretionary,NoHistoric},
  RawFeature={-calt;-tlig;+cv01;+cv32;+cv33;+cv34;+cv35;+cv36;+cv37}
]
\newCJKfontfamily\BookCodeCJK{Maple Mono NF CN}[
  Ligatures={NoCommon,NoContextual,NoDiscretionary,NoHistoric},
  RawFeature={-calt;-tlig;+cv01;+cv32;+cv33;+cv34;+cv35;+cv36;+cv37}
]
\makeatletter
\@namedef{FV@fontfamily@bookcode}{%
  \def\FV@FontScanPrep{}%
  \def\FV@FontFamily{\BookCodeLatin\BookCodeCJK}%
}
\makeatother
